% 第2章 2.4.3：C([0,1]) 中弱收敛加范数收敛不推出范数收敛。
% h(t)=max(0,2t-1)，g_n 为支撑在 [0,2/n]、峰值 1 位于 t=1/n 的尖峰；f_n=h+g_n。
% n≥4 时支撑不交，‖f_n‖_∞=1=‖h‖_∞，而 ‖f_n-h‖_∞=‖g_n‖_∞=1。图中 n=5（绿虚线）与 n=12（黄点划线）。
\begin{tikzpicture}[font=\figurefont\footnotesize]
\begin{axis}[width=76mm,height=40mm,axis lines=left,
  x axis line style={draw=FigureStroke,line width=.6pt},y axis line style={draw=FigureStroke,line width=.6pt},
  tick style={draw=FigureStroke},tick label style={font=\figurefont\scriptsize,text=FigureInk},
  xlabel={$t$},xlabel style={font=\figurefont\scriptsize,at={(axis description cs:1,0)},anchor=north west,yshift=1mm},
  xmin=0,xmax=1.02,ymin=0,ymax=1.15,xtick={0,0.5,1},ytick={0,0.5,1},
  ymajorgrids,grid style={draw=FigureGrid,line width=.4pt},
  legend style={draw=none,fill=none,font=\figurefont\scriptsize,cells={anchor=west},at={(1.02,0.5)},anchor=west}]
\addplot[draw=FigureBlue,line width=1.3pt] coordinates {(0,0) (0.5,0) (1,1)};
\addplot[draw=FigureGreen,line width=1.1pt,dashed] coordinates {(0,0) (0.2,1) (0.4,0) (0.5,0) (1,1)};
\addplot[draw=FigureYellow,line width=1.1pt,dash dot] coordinates {(0,0) (0.08333,1) (0.16667,0) (0.5,0) (1,1)};
\legend{$h$,$f_5=h+g_5$,$f_{12}=h+g_{12}$}
\end{axis}
\end{tikzpicture}
